
\SourceLeaf{452}\Article{3}
Diese sechs Bedingungen sind folgende:

Die beiden Achatplatten, auf deren obern Rändern die Zapfen der Nadel beim Beobachten zu liegen kommen, sollen durch die beiden Schraubenpaare, auf welche sie sich stützen, so regulirt sein, dass

1) ihre obern Ränder in Einer Ebene liegen,

2) dass diese Ebene normal gegen die Ebene des Verticalkreises ist, und

3) unterhalb des Mittelpunkts dieses Kreises liegt, mit einem der halben Zapfendicke gleichkommenden Abstande,

4) dass die Durchschnittslinie jener beiden Ebenen mit der Verticalachse einen rechten Winkel macht.

Es müssen ferner die Pfannen, vermittelst welcher man die Nadel von den Achatplatten abhebt und wieder auflegt, und die auf dem Heberahmen mit einiger Verschiebbarkeit aufgeschraubt sind. so regulirt sein, dass nach dem Auflegen der Nadel

5) ihre Achse normal gegen die zuletzt (in 4) genannte Durchschnittslinie wird (mithin in Verbindung mit der Bedingung 2 auch normal gegen die Ebene des Verticalkreises) und zugleich

6) den verticalen Durchmesser des Kreises trifft.

Die Bedingungen 1, 2, 4 zusammengenommen vertreten die Stelle der einen, dass bei genau senkrechter Stellung der aufrechten Drehungsachse eine horizontale Ebene die Ränder der beiden Achatplatten der Länge nach oder in zwei Linien berühren soll, insofern vorausgesetzt wird, dass die Ebene des Verticalkreises mit jener Drehungsachse parallel ist, also mit ihr zugleich senkrecht wird: man kann dies als die siebente Bedingung betrachten, welche man stillschweigend im Vertrauen auf die Geschicklichkeit des Künstlers vorauszusetzen pflegt, und zu deren Prüfung und, eventuell, Berichtigung das Instrument, wie es ist, keine Mittel darbietet.

\Article{4}
Bei den in diesem Aufsatze anzuführenden Beobachtungen war ich in Beziehung auf die Prüfung der angegebenen Bedingungen, in Ermangelung anderer Mittel, auf folgende Art zu Werke gegangen.

\SourceLeaf{453}Zur Prüfung der \textit{ersten} Bedingung gebrauchte ich das Planglas eines sogenannten künstlichen Horizonts, welches (nachdem vorher der Rahmen mit den Pfannen weggenommen war) so auf die Achatläger gelegt wurde, dass die mattgeschliffene Rückseite nach oben gekehrt war. Wenn die Bedingung nicht erfüllt ist, wird immer nur \textit{eine} Achatplatte nach der ganzen Länge, die andere am einen Endpunkte berührt werden, was man. wenn der Fehler nicht sehr gering ist, schon mit dem Auge erkennt; mehr Genauigkeit und Sicherheit gibt eine auf die Glasplatte gestellte Libelle, welche zeigt, ob diese zwei verschiedene Berührungslagen hat oder nur eine. Man sieht leicht, dass mit Hülfe dieser Libelle nach Erfüllung der ersten Bedingung auch die \textit{zweite} und \textit{vierte} geprüft werden kann.

Zur Prüfung der \textit{fünften} Bedingung muss die Nadel in zwei verschiedene Gleichgewichtsstellungen gebracht werden, und zwar solche, wo bei gleicher Lage der Zapfen auf den Lägern (oder indem dieselbe Nadelfläche vorne ist) die Nadel nur eine mässige Neigung gegen die Horizontallinie hat, aber das Ende, welches in der einen Lage auf der linken Seite war, bei der andern rechts zu stehen kommt. Man verschafft sich diese beiden Stellungen am bequemsten vermittelst angemessener Belastungen der Nadel, und erkennt das Erfülltsein der in Rede stehenden Bedingung daran, dass von der Schärfe jedes Nadelendes in der einen Stellung eben so viel vor den Rand des Kreises vortreten muss, wie in der andern. In Gegenden, wo nur eine mässige Inclination Statt findet, würden die betreffenden beiden Lagen schon durch blosse halbe Umdrehung des Instruments, so dass die Kreisfläche beidemal nahe am magnetischen Meridian ist, zu erhalten sein.

Eine ähnliche Prüfungsart lässt sich übrigens auch für die \textit{zweite} Bedingung anwenden, nur dass dabei zwei entgegengesetzte nahe verticale Stellungen der Nadel hergestellt sein müssen, wovon die eine sich von selbst ergibt, wenn man die Kreisebene nahe rechtwinklig gegen den magnetischen Meridian bringt, die andere entweder durch eine angemessene Belastung, oder durch Umkehren der Pole. Man sieht aber leicht, dass dieses Verfahren mit dem oben erwähnten vermittelst der Libelle nur dann gleichgeltend ist, wenn die siebente Bedingung erfüllt ist, und dass man also durch Verbindung beider Methoden eine Art von Prüfung dieser Bedingung selbst erhält, die freilich nur eine sehr unvollkommene sein kann, da sich das gleiche Vortreten der Nadelschärfe vor den Kreisrand nur schätzungsweise beurtheilen lässt.

\expandafter\def\csname HistoricalNoteText454\endcsname{Mit Ausnahme der vom 23. September.}%
\SourceLeaf{454}Dieselben combinirten Stellungen der Nadel dienen zugleich zur Prüfung der beiden übrigen Bedingungen; die \textit{sechste} Bedingung ist erfüllt, wenn jedes Nadelende in der ersten nahe horizontalen Stellung eben so weit von der innern Fläche des Kreises absteht, wie in der zweiten; für die \textit{dritte} Bedingung gilt ähnliches bei den nahe verticalen Stellungen. Offenbar würde zu der Prüfung hinreichen, die Abstände beider Nadelenden von der innern Kreisfläche unter sich bei Einer nahe horizontalen und Einer nahe verticalen Stellung zu vergleichen, wenn die beiden Nadelhälften genau gleich lang wären, aber bei unserm Instrumente, wo die Zwischenräume überhaupt so sehr klein sind, genügt dies nicht, und selbst eine sehr geringe Ungleichheit in den beiden Nadelhälften wird dabei schon bemerkbar.

\Article{5}
Wie schwer es ist, auf solche Art allen Bedingungen zugleich Genüge zu thun, erhellt schon aus dem Umstande, dass die zwei Schrauben, auf welchen jede Achatplatte ruht, nur acht Millimeter von einander abstehen, so dass, da die Weite eines Schraubengewindes $0.283$ Millimeter beträgt, schon eine halbe Umdrehung einer Schraube die betreffende Achatplatte um einen Grad wendet.

Sehr erleichtert wird aber das Geschäft durch eine eigne Vorrichtung, die ich erst später habe anfertigen lassen, und die dazu dient die Ränder der Achatplatten in Eine Ebene zu bringen und diese horizontal zu machen; ich halte mich aber jetzt nicht bei einer Beschreibung derselben auf, da sie für die gegenwärtigen Beobachtungen\HistoricalNote{454} noch nicht hatte benutzt werden können. Eine zweite gleichfalls erst nach dem Schluss der Beobachtungen fertig gewordene Vorrichtung dient zu einer scharfen Bestimmung der Abweichung des Hauptkreises von der verticalen Lage. Sie hat diese Abweichung zu \textit{zehn Minuten} ergeben, aber die Wegschaffung der Abweichung wird erst eine Abänderung am Instrumente erfordern. Übrigens kann der Einfluss dieser Abweichung auf die Inclinationen nicht einmal eine Secunde betragen.

Überhaupt darf ich nicht unbemerkt lassen, dass kleine Fehler in den verschiedenen Berichtigungen nur einen kaum merklichen Einfluss auf die Inclina%
\SourceLeaf{455}%
tionsbestimmungen haben können. Der Einfluss, welchen auf die Stellung der Nadel ein Theil der Fehler hat, ist in Beziehung auf diese nur eine Grösse der zweiten Ordnung, und die Wirkung der andern, namentlich einer Excentricität, und einer Neigung der die Achatplatten berührenden Ebene in dem Sinn parallel mit der Ebene des Kreises (Fehler gegen die Bedingungen 3, 6 und 4) werden durch die Combination der einzelnen Beobachtungsstücke völlig eliminirt. Ich kann daher dem Urtheil \textsc{Horner}'s, dass \textit{vor allem} auf die Wegschaffung dieses letzten Fehlers zu sehen sei (Physik. Wörterb. 5. Band, S. 759) nicht beistimmen, sondern betrachte diesen Fehler als denjenigen, an dessen vollkommener Wegschaffung am wenigsten gelegen ist.

\Article{6}
Die hier aufzuführenden Inclinationsbeobachtungen sind sämmtlich im Freien an dem in den \textit{Resultaten} 1840. II. [S. 433 d. B.] bezeichneten Platze angestellt; ein Schirmdach hielt die Sonnenstrahlen von dem Instrumente ab. Dieses wurde auf dem Steine so aufgestellt, dass die gerade Linie durch zwei Fussspitzen nahe senkrecht gegen den magnetischen Meridian wurde, für welche Stellung die Plätze der drei Füsse bezeichnet waren. Die genaue magnetische Orientirung des Instruments wurde durch eine demselben beigegebene Hülfsnadel erhalten, die mit einem Achathütchen auf eine Spitze aufgehängt wird; der Träger dieser Spitze hat zwei kurze cylindrische Seitenarme, die in die beiden Pfannen eingelegt werden, wodurch sich die Spitze in Folge des Gewichts des frei herabhängenden Theils des Trägers von selbst vertical stellt. Ich habe öfters mit dieser Orientirungsart auch die sonst übliche durch correspondirende Neigungen in zwei nahe gegen den magnetischen Meridian senkrechten Stellungen des Verticalkreises verbunden und immer nur ganz unerhebliche Unterschiede gefunden, woraus hervorgeht, dass die Hülfsnadel hinlänglich empfindlich ist und keine constante Abweichung hervorbringt. Eine geringe Abweichung der Verticalebene, in welcher man beobachtet, von dem ohnehin während der Beobachtungen nicht ganz unveränderlichen magnetischen Meridian hat übrigens auf die Neigung der Inclinationsnadel nur einen als ganz unmerklich zu betrachtenden Einfluss von der zweiten Ordnung.

\SourceLeaf{456}\Article{7}
Das Zusammenfallen des Schwerpunkts einer Nadel mit der Drehungsachse können die geschicktesten Künstler nur näherungsweise bewirken: es bleibt fast immer eine Abweichung zurück, deren Einfluss auf die Einstellung der Nadel durch die Combination von Beobachtungen unter mehrfach gewechselten Umständen ermittelt oder eliminirt werden soll: zu diesen abgeänderten Umständen gehört wesentlich die Umkehrung der Pole der Nadel. Unter sonst gleichen Umständen ist jener Einfluss desto stärker, je schwächer die Nadel magnetisirt ist; da man aber nicht befugt ist, anzunehmen. dass die Stärke des Nadelmagnetismus nach dem Umkehren der Pole wieder eben so gross wird, wie vorher, so ist eine genaue Reduction der Beobachtungen von der Kenntniss des Verhältnisses dieser Stücke abhängig. Man gelangt dazu durch Beobachtung der Schwingungsdauer der Nadel: ich habe aber aus mehrern Gründen \textit{horizontalen} Schwingungen den Vorzug gegeben, und zu deren Beobachtung einen besondern von Hrn. Inspector \textsc{Meyerstein} verfertigten Apparat angewandt. Die Nadel schwingt in einem hölzernen Kasten mit verglasten Deckeln, und liegt dabei auf einem leicht gearbeiteten Bügel, der an einem $270$ Millimeter langen von einer Glasröhre gegen Luftzug geschützten Seidenfaden hängt, und ihre Enden spielen während der Schwingungen an zwei Gradbogen, deren jeder $40$ Grad umfasst, in halbe Grade getheilt ist, und $5$ Minuten mit Sicherheit zu schätzen verstattet. Die Schwingungsdauer jeder Nadel wurde vor und nach dem Umstreichen jedesmal aus $150$ in drei Sätze vertheilten Schwingungen bestimmt, die nach gehöriger Reduction auf unendlich kleine Bögen stets vortrefflich übereinstimmende Resultate geben. Angefangen wurde gewöhnlich mit einem Schwingungsbogen von etwa $36$ Grad, und es verdient hier wohl bemerkt zu werden, dass, im Gegensatz gegen die in den \textit{Resultaten} 1837. IV [S. 385 d. B.] erwähnten Erfahrungen an schwereren Stäben, die Abnahme des Schwingungsbogens an allen Tagen und Nadeln mit fast gleicher Geschwindigkeit erfolgte, so dass die Zeit, innerhalb welcher der Bogen auf seinen vierten Theil herabkam, mit geringen Schwankungen $14$ Minuten betrug. Übrigens wurden diese Schwingungsbeobachtungen immer in der Sternwarte auf einem Steinpostamente angestellt, indem es dabei nicht sowohl auf die absolute Dauer, als auf das Verhältniss ankommt, welches von den kleinen in diesem Local möglicherweise Statt findenden fremden Einflüssen nicht merklich afficirt werden kann.

\expandafter\def\csname HistoricalNoteText457\endcsname{Das auffallendste Beispiel dieser Art wird in dem \textit{Fifth Report of the British association for the advancement of Science} S. 142 angeführt, wo acht Nadeln, mit welchen Capitaine Ross in London die Inclination bestimmte, Unterschiede bis zu 41 Minuten ergaben, obgleich die Beobachtungen mit jeder einzelnen Nadel zahlreich und unter sich gut übereinstimmend waren. Die Ursache dieser sonderbaren Erscheinung, über welche näheres Detail nicht mitgetheilt ist, hat man in England der nicht vollkommen cylindrischen Gestalt der Zapfen beigemessen, und gerade deshalb drehbare Zapfen versucht.}%
\SourceLeaf{457}\Article{8}
Bei den im Sommer 1842 angestellten Inclinationsbeobachtungen be\-zweckte ich ausser der Feststellung der für diese Zeit geltenden magnetischen Inclination zugleich die Bestimmung des Grades der Genauigkeit, welche mit dem angewandten Instrument erreicht wird. Es schien mir nicht genügend, die Zuverlässigkeit der Endresultate, auf welche so mancherlei Umstände Einfluss haben, nach den Unterschieden abzuschätzen, die sich in den Einstellungen der Nadel bei wiederholtem Abheben \textit{vermittelst des Pfannenhebels} ergeben; eben so wenig aber kann zu diesem Zwecke die blosse Vergleichung der Resultate dienen, die man für die Inclination aus den Beobachtungen verschiedener Tage erhält, da sich dabei die zufälligen dem Instrument beizumessenden Beobachtungsfehler mit den wirklichen Schwankungen der Inclination selbst vermischen. Ich war ferner begierig zu erfahren, ob meine vier Nadeln übereinstimmende, oder wie es einigen Beobachtern begegnet ist\HistoricalNote{457}, entschieden und bedeutend ungleiche Resultate geben würden.

Diese Rücksichten haben mich bewogen, eine von der gewöhnlichen etwas abweichende Anordnung der Beobachtungen zu wählen; das Wesentliche des Unterschiedes ergibt sich aus folgendem.

Gewöhnlich beobachtet man den Stand der Nadel, d. i. die Stellung beider Spitzen gegen die Theilung des Kreises, in vier verschiedenen Combinationen der Stellung des Kreises und der Art des Einliegens der Nadel, indem die getheilte Fläche des erstern und die gezeichnete Fläche der letztern nach Osten oder Westen, nach gleicher oder nach entgegengesetzter Weltgegend gekehrt sein können. Dieselben Combinationen werden nach dem Umkehren der Pole wiederholt, so dass zusammen $16$ Ablesungszahlen vorliegen, aus welchen man, in so fern sie nicht in Folge einer starken Abweichung des Schwerpunkts der Nadel von ihrer Zapfenachse grosse Verschiedenheiten darbieten, das einfache arithmetische Mittel für die Inclination annimmt, oder im entgegengesetzten Falle eine künstlichere Rechnung anwendet. Es versteht sich, dass jede der $16$ Zahlen
\SourceLeaf{458}selbst schon das Mittel aus einer kleinern oder grössern Anzahl von Einstellungen sein kann, die man in jeder Combination durch wiederholtes Aufheben erhält.

Hievon unterscheidet sich das von mir befolgte Verfahren dadurch, dass ich an jedem Tage mit zwei Nadeln beobachtet habe, ohne zwischen den Beobachtungen die Pole umzukehren; das Umkehren der Pole geschah zwischen zwei auf einander folgenden Beobachtungen und zwar wechselsweise immer nur an einer Nadel. Man sieht, dass auf diese Art die Beobachtungen von vier Tagen alle Combinationen der verschiedenen Polarisirungen beider Nadeln umfassen, wie dies mit den Nadeln 1 und 3 vom 6. bis 9. Julius, und mit den Nadeln 2 und 4 vom 17. bis 20. Julius geschehen ist. Eine Fortsetzung ähnlich combinirter Abwechslungen durch acht Beobachtungstage, wie mit den Nadeln 1 und 2 vom 20. Mai bis 5 Junius, und mit den Nadeln 3 und 4 vom 8. bis 25. Junius ausgeführt ist, gab also jede Combination der Polarisirungen zweimal. Die Beobachtungen an jedem Tage wurden so geordnet, dass die Resultate aus beiden Nadeln, so viel thunlich, gleichzeitig wurden. Dies wurde dadurch erreicht, dass zuerst die oben erwähnten vier Combinationen an der einen Nadel durchbeobachtet wurden, und zwar jede mit viermal wiederholter Auflegung; sodann die ähnlichen Combinationen an der zweiten Nadel unter achtmal wiederholter Auflegung; endlich wiederum an der ersten Nadel dieselben Combinationen, aber in verkehrter Ordnung und unter viermal wiederholter Auflegung.

Bei dieser Einrichtung geben die Beobachtungen Eines Tages für sich allein noch keine Inclinationsbestimmung; allein wenn damit die Beobachtungen des folgenden Tages verbunden werden, so lässt offenbar die nicht umgestrichene Nadel erkennen, um wie viel die Inclination an den beiden Tagen ungleich war, und die einseitigen Beobachtungen an der andern können danach auf Einen Zeitpunkt reducirt, und also vollständig gemacht werden. Zu einer strengern die Gesammtheit der Beobachtungen von allen 24 Tagen umfassenden Behandlung wird aber erst das gegenseitige Verhalten der partiellen Beobachtungsresultate näher erörtert werden müssen.

\Article{9}
Diese in mehr als einer Beziehung wichtige Entwickelung wird sich am bequemsten an ein Beispiel anknüpfen lassen, entnommen von einer auf die ge%
\SourceLeaf{459}%
wöhnliche Art angestellten Beobachtung, dergleichen von mir auch an mehrern Tagen gemacht sind.

Ich wähle dazu die Beobachtung mit der Nadel 1 vom 23. September 1842 Vormittags von $8\frac{1}{2}$ bis $11$ Uhr. Die magnetische Orientirung wurde auf die im 6. Art. angezeigte Art mit der Hülfsnadel erhalten, und der Index des Azimuthalkreises (dessen von der Linken nach der Rechten wachsende Grade ich, wie schon oben bemerkt ist, von $0$ bis $360^0$ durchzähle) zeigte bei der Stellung des Verticalkreises im Meridian, die getheilte Seite nach Osten gekehrt, $90^0\ 5'$.

Ausser den gewöhnlichen acht Combinationen im magnetischen Meridian machte ich an diesem Tage noch eben so viele in der gegen denselben normalen Verticalebene: ich nehme diese Beobachtungen hier mit auf, da sie zu mehrern Erörterungen Gelegenheit geben. Die Nadel ist (eben so wie die drei andern) auf einer Seite mit den Buchstaben $A$, $B$ an den Enden gezeichnet, wodurch die Polarisirung und Einlegungsart bequem unterschieden werden kann. In jeder der $16$ Combinationen wurde die Nadel fünfmal mit dem Pfannenhebel auf die Achatplatten gelegt: in der folgenden Übersicht gebe ich aber nur die Mittelwerthe aus den zusammengehörigen Einstellungen.

\begin{center}
\textit{Nadelende $B$ Nordpol.}\par\smallskip
{\fontsize{9}{11}\selectfont
\begin{tabular}{c||cc||cc}
Azimuthal&\multicolumn{4}{c}{Bezeichnete Nadelfläche}\\
Kreis&\multicolumn{2}{c||}{vorne}&\multicolumn{2}{c}{hinten}\\
&unten&oben&unten&oben\\\hline
$90^0\ 5'$&$67^0\ 27'\ 54''$&$67^0\ 29'\ 36''$&$67^0\ 45'\ 39''$&$67^0\ 44'\ 51''$\\
$180\ 5$&$89\ 52\ 39$&$89\ 52\ 51$&$90\ 12\ 30$&$90\ 10\ 30$\\
$270\ 5$&$112\ 18\ 39$&$112\ 16\ 45$&$112\ 38\ 51$&$112.\ 33\ 54$\\
$0\ 5$&$89\ 58\ 33$&$89\ 57\ 48$&$90\ 13\ 27$&$90\ 10\ 54$
\end{tabular}}
\par\medskip
\textit{Nadelende $A$ Nordpol.}\par\smallskip
{\fontsize{9}{11}\selectfont
\begin{tabular}{c||cc||cc}
$90^0\ 5'$&$68^0\ 2'\ 51''$&$68^0\ 2'\ 33''$&$67^0\ 35'\ 15''$&$67^0\ 37'\ 0''$\\
$180\ 5$&$90\ 14\ 48$&$90\ 12\ 21$&$89\ 51\ 12$&$89\ 51\ 36$\\
$270\ 5$&$112\ 27\ 21$&$112\ 22\ 33$&$112\ 7\ 6$&$112\ 5\ 33$\\
$0\ 5$&$90\ 16\ 15$&$90\ 14\ 0$&$89\ 53\ 54$&$89\ 54\ 18$
\end{tabular}}
\end{center}

\Needspace{6\baselineskip}
\SourceLeaf{460}Die Dauer einer horizontalen Schwingung wurde gefunden
\begin{center}
\begin{tabular}{lll}
vor den Beobachtungen&$.\quad.\quad.\quad.\quad.\quad.\quad.\quad.\quad.$&$5''83555$\\
nach den Beobachtungen&$.\quad.\quad.\quad.\quad.\quad.\quad.\quad.\quad.$&$5.87416$
\end{tabular}
\end{center}

\Article{10}
Ich verweile nun zuerst bei den Unterschieden zwischen den Ablesungen der untern und obern Spitze, welche davon abhängen, dass die Zapfenachse weder durch den Mittelpunkt der Theilung, noch durch die die beiden Spitzen der Nadel verbindende gerade Linie geht. Bezeichnen wir mit $x$, $y$ die Coordinaten des Schnittes der Zapfenachse mit der Kreisebene relativ gegen den Mittelpunkt der Theilung, ausgedrückt in Bogentheilen der innern Kreisperipherie, und zwar $x$ parallel mit dem Diameter durch die beiden Nullpunkte und positiv nach der rechten, $y$ parallel mit dem Diameter durch die beiden $90^0$ Punkte und positiv nach oben; ferner mit $180^0-z$ den Winkel zwischen den beiden durch die Zapfenachse und die Spitzen $A$ und $B$ gelegten Ebenen, so verstanden, dass, indem man sich die Nadel horizontal und die gezeichnete Seite nach oben gekehrt denkt, in dem Sinne von der linken nach der rechten von $A$ nach $B$ gezählt wird; endlich mit $l$ das Mittel zwischen den beiden Ablesungen: so wird der Unterschied derselben (so verstanden, dass die untere Ablesung von der obern abgezogen wird)
\[
=2x\sin l+2y\cos l\pm z
\]
wo das obere Zeichen gilt, wenn zugleich die gezeichnete Fläche vorne und $A$ oben (also hier Südpol), oder jene hinten und $B$ oben ist, das untere Zeichen in den beiden andern Fällen.

Die obigen Beobachtungen geben so $16$ Gleichungen, aus welchen nach der Methode der kleinsten Quadrate gefunden wird
\[
\begin{aligned}
x&=-38''3\\
y&=+153.2\\
z&=+75.4
\end{aligned}
\]
\Needspace{20\baselineskip}
Die Vergleichung gibt dann, wenn man nach der Grösse von $l$ ordnet,
\SourceLeaf{461}\begin{center}
{\fontsize{9}{11}\selectfont
\begin{tabular}{c|r|r|r}
$l$&Beobachtung&Rechnung&Fehler\\\hline
$67^0\ 28'\ 45''$&$+102''$&$+122''$&$-20''$\\
$67\ 36\ 7$&$+105$&$+121$&$-16$\\
$67\ 45\ 15$&$-48$&$-30$&$-18$\\
$68\ 2\ 42$&$-18$&$-32$&$+14$\\
$89\ 51\ 24$&$+24$&$0$&$+24$\\
$89\ 52\ 45$&$+12$&$-1$&$+13$\\
$89\ 54\ 6$&$+24$&$-1$&$+25$\\
$89\ 58\ 10$&$-45$&$-1$&$-44$\\
$90\ 11\ 30$&$-120$&$-153$&$+33$\\
$90\ 12\ 10$&$-153$&$-153$&$0$\\
$90\ 13\ 34$&$-147$&$-153$&$+6$\\
$90\ 15\ 7$&$-135$&$-153$&$+18$\\
$112\ 6\ 19$&$-93$&$-111$&$+18$\\
$112\ 17\ 42$&$-114$&$-112$&$-2$\\
$112\ 24\ 57$&$-288$&$-264$&$-24$\\
$112\ 36\ 22$&$-297$&$-264$&$-33$
\end{tabular}}
\end{center}

Die Summe der Quadrate der übrig bleibenden Fehler ist $7924$, woraus man schliesst, dass der mittlere Fehler der Differenz zweier Mittel aus fünf Ablesungen
\[
=\BareRoot{\frac{7924}{13}}=24''7
\]
und der mittlere Fehler der einfachen Ablesung Einer Spitze
\[
=\BareRoot{\frac{5}{2}.\frac{7924}{13}}=39''0
\]
angenommen werden kann, eine in der That sehr befriedigende Genauigkeit, welche durch ähnliche Discussion der Beobachtungen von andern Tagen nicht nur bestätigt, sondern zuweilen noch übertroffen wird. Es mag jedoch dabei bemerkt werden, dass die Erreichung einer solchen Übereinstimmung wesentlich von dem Umstande abhängt, dass das Abheben der Nadel immer nur dann geschieht, wenn sie in Ruhe oder ihrer Ruhestellung nahe ist. Ohne diese Vorsicht würde die Nadel, deren Schwingung in einem Rollen der Zapfen auf dem Lager besteht, an einer andern Stelle des Lagers, als wo sie niedergelegt wird, zur Ruhe kommen, und also das Excentricitätselement $x$ ein veränderliches sein.

Man erhält auf die hier angegebene Art allerdings die Werthe der Excentricitätselemente $x$ und $y$ mit vieler Genauigkeit, allein diese Werthe können
